\documentclass{article}
\usepackage[utf8]{inputenc}
\usepackage[T1]{fontenc}
\usepackage{geometry}
\geometry{
    a4paper,
    left=2.5cm,
    right=2.5cm,
    top=2.5cm,
    bottom=2.5cm
}
\usepackage{graphicx}
\usepackage{float}
\usepackage{hyperref}
\usepackage[english]{babel}
\usepackage{amsmath}
\usepackage{amssymb}
\usepackage{array}
\usepackage{booktabs}
\usepackage{caption}
\usepackage{subcaption}
\usepackage{enumitem}
\usepackage{fancyhdr}
\usepackage{titlesec}
\usepackage{url}
\usepackage{textcomp}
\usepackage{xcolor}
\usepackage{listings}
\usepackage{multirow}
\usepackage{longtable}
\usepackage{tabularx}
\usepackage{times}

% Set up page style
\pagestyle{fancy}
\fancyhf{}
\rfoot{\thepage}
\renewcommand{\headrulewidth}{0pt}

% Title formatting
\titleformat{\chapter}[display]{\normalfont\bfseries}{\chaptertitle\thechapter}{0pt}{\Huge}
\titlespacing*{\chapter}{0pt}{50pt}{40pt}

% Document info
\title{Template Document}
\author{Academic Support Agent}
\date{\today}

\begin{document}
\maketitle

% El resumen debe ser lo  ltimo en aparecer antes de empezar con el contenido del art culo.
\begin{abstract}
Este archivo sirve como plantilla para art culos que se quieran presentar a la revista TEMat de la ANEM.
Presentamos de manera simple c mo se debe usar el formato de la revista, as  como las instrucciones habituales y todo lo relacionado con la redacci n de art culos en \LaTeX.
Se puede utilizar la propia plantilla como base para trabajar.
Tambi n discutimos los paquetes necesarios, de los cuales la mayor a se cargan autom ticamente, y algunos paquetes prohibidos.
Se dan varios ejemplos de uso de la mayor parte de los elementos presentados.

Para asegurar un estilo uniforme en todos los art culos de la revista, el uso de la plantilla es obligatorio, y est  totalmente prohibido cambiar los ajustes de espaciado y estructura.
\end{abstract}
\begin{abstract}[english]
This file serves as a template for articles that are planned to be presented to ANEM's TEMat journal.
We present how to use the journal's format in a simple way, as well as the usual instructions and everything related to writing articles using \LaTeX.
You might use this template as a basis for your work.
We also discuss about the packages that are needed, considering that most of them are automatically loaded, and some forbidden packages.
We show many use-case examples for most of the presented elements.

To ensure a uniform style in all articles in the journal, using the template is compulsory and it is completely forbidden to change the spacing and structure settings.
\end{abstract}
\maketitle

\section{Introducci n}

Esta es la plantilla de art culos para la revista TEMat.
En ella recogemos todas las instrucciones b sicas sobre c mo crear art culos en \LaTeX{} de acuerdo a los criterios editoriales de TEMat.
Hay varios documentos que forman parte de la plantilla.
Los m s importantes son el documento \verb+TEMat-article.cls+, que forma una clase que incluye todo el formato de la revista y \textbf{NO} se debe modificar; el documento \verb+PlantillaTEMat.tex+, que sirve para crear art culos, y el documento \verb+plantilla.bib+, que sirve para crear bibliograf as que despu s se importan al art culo.
Adem s, est n los archivos \verb+latexmkrc+ (necesario para los que utilicen \verb+latexmk+), y \verb+triangulo.tex+ y \verb+triangulo-ps.tex+ (que proporcionan ejemplos de figuras en \verb+tikz+ y en \verb+pstrics+), as  como los diagramas \verb+triangulo-precompilado.pdf+ y \verb+triangulo-ps-precompilado.pdf+ como muestra para la inclusi n de figuras.
Finalmente, el documento \verb+PlantillaTEMat.pdf+ permite ver el resultado despu s de compilar.
Si no puedes compilar la plantilla, eso quiere decir que faltan paquetes en tu ordenador, y se deben instalar.
Si el resultado de la compilaci n es distinto del ofrecido en el pdf, tambi n hay alg n problema.
Se puede compilar con \verb+pdflatex+, \verb+lualatex+ o \verb+xelatex+ (preferimos uno de los dos primeros, y en principio nosotros compilaremos con el segundo).
No damos soporte a compilaciones que no sean en modo pdf (\verb|latex+dvips+ps2pdf| u otras).

Esta plantilla (especialmente el documento \verb+TEMat-article.cls+) podr a ser revisada y modificada a lo largo del tiempo.
Recomendamos a los autores que descarguen la versi n m s reciente de la clase para trabajar.
En caso de encontrar un error o tener alg n problema, contactad con los editores de TEMat a trav s del correo \email{publicaciones@anemat.com}.

\section{Portada del art culo}

Cada art culo tendr  su propia portada (para que se puedan descargar por separado y quede recogida toda la informaci n del art culo en dicha portada), parecida a la de esta plantilla.
En esta p gina deber  aparecer toda la informaci n del autor o los autores y del art culo en cuesti n.
N tese que toda la informaci n se puede introducir antes o despu s del \verb+\begin{document}+, salvo los res menes que deben ir siempre despu s.
La portada se genera despu s de los res menes mediante el uso del comando \verb+\maketitle+.

Se debe poner el t tulo del art culo con el comando \verb+\title{}+.
Si se trata de un t tulo largo (tanto que no cabe en la cabecera de las p ginas de la revista), se deber  proporcionar tambi n un t tulo corto.
Para ello, se deben rellenar los dos campos como sigue: \verb+\title[T tulo corto, opcional]{T tulo del art culo}+.
Se puede ver un ejemplo de esto y de todo lo posterior en el archivo \verb+PlantillaTEMat.tex+ de esta plantilla.

Es necesario incluir el nombre y los apellidos para cada autor, as  como su afiliaci n (en qu  universidad (u otro lugar) estudia o trabaja).
Tambi n existe la posibilidad de incluir un correo electr nico, por si alguien quiere contactar al autor para preguntar m s sobre el tema, pero poner esto es opcional.
Para introducir estos datos, se deben utilizar los siguientes comandos:
\begin{itemize}
\item \verb+\author{Apellido o Apellidos, Nombre}+ para el nombre (este comando debe ser el primero en aparecer para cada autor);
\item \verb+\affiliation{Afiliaci n}+ para indicar la afiliaci n acad mica o lugar de trabajo;
\item \verb+\email{email}+, de forma opcional, para indicar una direcci n de correo electr nico de contacto, y
\item \verb+\authornote{Nota al pie}+, de forma opcional, para a adir una breve nota al pie sobre esa persona.
\end{itemize}
En el mundo de las matem ticas, la pr ctica m s habitual es poner a los autores por orden alfab tico de apellido, sin darle m s  nfasis a uno que a los otros, y nosotros aconsejamos esta pr ctica; sin embargo, mantendremos el orden en que est n en el env o.

Cuando solamente hay un autor, debe usarse el comando \verb+\author*+ (con asterisco) para indicar su nombre. Si hay m s de uno, la persona de contacto deber  usar el comando con asterisco y el resto deber n usarlo sin asterisco.

Deber  haber una lista con las palabras clave del art culo, que se debe definir con \verb+\keywords{}+, separ ndolas por comas (y sin a adir punto final).
Tambi n se debe incluir la clasificaci n del tema del art culo seg n la MSC2010 (2010 Mathematics Subject Classification) usando \verb+\msc{}+.
Para los que no conozcan esto, se debe buscar en la p gina \url{http://www.ams.org/msc/msc2010.html} el c digo (o c digos) que mejor se correspondan con la tem tica del art culo que se presenta.
De ser posible, se debe poner uno solo; en caso de haber varios temas relacionados, recomendamos que no se pongan m s de cinco c digos.

Debajo de esta informaci n aparecen la fecha de recepci n y de aceptaci n del art culo presentado.
Las que hay en la plantilla est n a modo de ejemplo; estos campos los rellenar n los editores de la revista para cada art culo.
En la versi n final de los art culos, aparece tambi n la informaci n bibliogr fica del trabajo.

En la misma p gina aparecer n los agradecimientos (en caso de haberlos), as  como las fuentes de financiaci n (en caso de haberlas).
Para introducirlas, se utiliza \verb+\acknowledgements{}+.

Finalmente, se debe proporcionar un resumen (o \textit{abstract}) del art culo que se env a. Este resumen debe redactarse tanto en castellano como en ingl s.
Para ello, se define un entorno de resumen delimitado por \verb+\begin{abstract}+ y \verb+\end{abstract}+ para el idioma principal (castellano por defecto) y \verb+\begin{abstract}[english]+ \verb+\end{abstract}+ para ingl s.
El resumen debe aportar informaci n suficiente como para saber de qu  trata el art culo y cu les son las aportaciones principales, pero debe ser lo suficientemente corto como para que no se salga de la portada ninguno de los campos que deben estar en ella.

Una vez se han indicado todos los campos de la portada, debe a adirse el comando \verb+\maketitle+, que aparecer  despu s de todos ellos y justo antes de empezar con el contenido del art culo.

Para anonimizar el documento para la revisi n, basta con cargar la clase \verb+TEMat-article+ con la opci n \verb+anonimo+ (aunque tambi n \verb+anon+, \verb+anonymous+, \verb+revision+ o \verb+review+).
Es decir, habr a que poner al principio del documento el siguiente comando: \verb+\documentclass[anon]{TEMat-article}+

Insistimos: hay ejemplos de todo en \verb+PlantillaTEMat.tex+.



\section{Creando secciones...}

Todo art culo debe seguir una estructura que convierta su lectura en algo agradable y permita seguir los argumentos con mayor facilidad.
Habitualmente, esta estructura se ve reforzada utilizando secciones y subsecciones.
Para crear una secci n basta con usar el comando \verb+\section{nombre}+.
Para crear subsecciones, se debe usar \verb+\subsection{nombre}+.
Existe la posibilidad de llegar hasta las subsubsecciones, pero recomendamos encarecidamente que no se haga esto.

\subsection{...y subsecciones...}

Obviamente, en cada secci n se puede introducir tanto texto como quer is.

\subsection{...y m s subsecciones}

\section{Matem ticas}

La principal herramienta para un matem tico al usar \LaTeX{} son las f rmulas matem ticas.
Todas las f rmulas aparecer n autom ticamente en letra it lica, distinguiendo as  el entorno de matem ticas del texto normal.
Para a adir f rmulas, basta con usar \verb+\(+ y \verb+\)+, y escribir la f rmula matem tica entre estos delimitadores.
En su defecto, se pueden usar tambi n \verb+$+ y \verb+$+.
As , se pueden escribir peque as f rmulas en entorno matem tico, como \(x^3+y^3=z^3\) o $\nume^{\iunit\pi}=-1$.
Si se quiere poner f rmulas centradas y separadas del resto del p rrafo, se deben emplear los delimitadores \verb+\[+ y \verb+\]+ (\textbf{no} puedes usar \verb+$$+ y \verb+$$+, puede estropear el formato de la revista\footnote{En realidad, no quieres usar los dobles d lares en ning n documento.  El motivo? A diferencia de otros entornos, no tienen bien ajustado el espaciado.}), para conseguir formas del estilo de
\[\sum_{i=1}^n\frac{1}{i}\approx\log n.\]
S , las f rmulas son parte del texto, de modo que van seguidas de los signos de puntuaci n correspondientes.
Si adem s se quiere que la f rmula quede numerada, se puede usar \verb+\begin{equation}+ y \verb+\end{equation}+, de modo que resulte
\begin{equation}\label{equa:ejemplo1}
\int_{1}^x\frac{1}{t} \diff t = \log x.
\end{equation}
Existen otras muchas formas m s complejas de introducir ecuaciones, con entornos como \verb+align+ y otros muchos.
Recomendamos buscar informaci n sobre estos entornos por internet.
Tambi n se debe buscar informaci n sobre c mo introducir caracteres especiales y diferentes s mbolos matem ticos.
Por ejemplo, para referirnos a los n meros naturales, enteros, etc. ($\mathbb{N}$, $\mathbb{Z}$, $\mathbb{Q}$...) se utiliza el comando \verb+\mathbb{}+.
Se puede encontrar informaci n sobre estos s mbolos en p ginas como \url{https://www.artofproblemsolving.com/wiki/index.php/LaTeX:Symbols} o \url{http://detexify.kirelabs.org/classify.html} (que permite dibujar el s mbolo que buscas).

A menudo es frecuente querer referirse a una ecuaci n que se ha escrito antes (para eso se utiliza la numeraci n de las f rmulas).
As , para referirnos a la f rmula \eqref{equa:ejemplo1}, tenemos que escribir el comando \verb+\eqref{formula}+, donde el nombre de la f rmula se asigna con un \verb+\label+, como pod is comprobar en la ecuaci n anterior en el archivo \verb+PlantillaTEMat.tex+.
Alternativamente, se puede utilizar \verb+\cref{nombre}+.

\subsection{Entornos matem ticos}

Adem s del entorno b sico para escribir matem ticas, los matem ticos siempre escriben teoremas, corolarios, demostraciones..., y para cada uno de los tipos de enunciado hay un entorno propio.
\textbf{Se deben usar estos entornos para mantener el estilo de la revista}.
Para usar cada uno de ellos se debe empezar con un \verb+\begin{nombreentorno}+ y terminar con un \verb+\end{nombreentorno}+.
Por ejemplo, podemos escribir teoremas como los que siguen.

\begin{teorema}[del punto gordo]\label{teor:ejemplo1}
Dadas tres \emph{rectas} $r$,\/ $s$ y\/ $t$ en\/ $\mathbb{R}^2$, existe un punto tal que las tres rectas se intersecan en  l (vaya, que se cortan).
\end{teorema}

\begin{demostracion}
Basta con tomar un punto suficientemente gordo.
\end{demostracion}

\begin{teorema}\label{teor:ejemplo2}
Contrario a la creencia com n,
\[\frac{0}{0}=2.\]
\end{teorema}

\begin{demostracion}[Demostraci n por simplificaci n de ceros]
Lo demostramos con la siguiente cadena de igualdades:
\[\frac{0}{0}=\frac{100-100}{100-100}=\frac{(10+10)(10-10)}{10\,(10-10)}=\frac{20}{10}=2.\qedhere\]
\end{demostracion}

\begin{nota*}
En general, vamos a intentar que teoremas de este estilo que mand is no se publiquen. Os vigilamos.
\end{nota*}

Para a adir informaci n al encabezado del teorema (como un nombre, autor, a o...) se debe incluir entre corchetes (\verb+[]+) despu s del \verb+\begin{teorema}+, como en el \autoref{teor:ejemplo1}.
Adem s, hay un detalle importante en cuanto a las demostraciones.
El cuadradito negro aparecer  autom ticamente al final de la demostraci n, pero no estar  bien colocado si la demostraci n acaba con una f rmula.
En ese caso, se debe incluir la instrucci n \verb+\qedhere+ al final de la  ltima f rmula, como sucede en la demostraci n del \autoref{teor:ejemplo2}.

Los entornos que hemos habilitado son teorema, proposici n, corolario, lema, propiedad, conjetura, hecho, criterio, axioma, definici n, ejemplo, ejercicio, problema, pregunta, hip tesis, observaci n, convenci n, nota, notaci n, afirmaci n, caso, algoritmo, soluci n y, por supuesto, demostraci n.
Para crear cualquiera de ellos, pon su nombre (sin tilde) en \verb+nombreentorno+, o sus nombres habituales en ingl s, o alguna de las abreviaciones m s habituales (no garantizamos que tu abreviaci n favorita vaya a funcionar; si no funciona, escribe el nombre completo).
No, no hay m s entornos.
Si crees que alguno m s es absolutamente necesario, ponte en contacto con nosotros.

En general, cada vez que se empiece a trabajar en un entorno, este aparecer  numerado.
Si el autor quiere que no sea as , basta con a adir un asterisco (\verb+*+) despu s de \verb+nombreentorno+ (por ejemplo, \verb+teorema*+).

Tambi n es muy frecuente querer referirnos a resultados que ya hemos escrito, como el \autoref{teor:ejemplo1}. Para hacerlo, basta con usar \verb+\autoref{nombreteorema}+ o \verb+\cref{nombreteorema}+, donde al teorema se le ha dado un nombre usando \verb+\label+ (de nuevo, puedes comprobar c mo en el archivo \verb+PlantillaTEMat.tex+).
Haciendo esto, si tienes el documento en pdf, puedes pinchar en \autoref{teor:ejemplo2} para que te muestre directamente el resultado al que se refiere, cosa muy  til a la hora de leer art culos.

\subsection{Sistema Internacional}

Para escribir n meros grandes o magnitudes f sicas, hemos cargado el paquete \verb+siunitx+.
Si quer is saber c mo se usa, solamente ten is que echarle un vistazo a la documentaci n~\cite{siunitx}.
Brevemente, para escribir un n mero como \num{10000}, usa \verb+\num{10000}+ (y as  se ver  el espacio recomendado cuando hay m s de cuatro cifras).
Esperamos que el \SI{100}{\percent} de los autores que usen unidades en sus art culos aprovechen su uso.
Especialmente porque, como dice \citet{bezosOM}, las unidades se escriben en letra redonda y con un espacio entre el n mero y la unidad.
Por ejemplo, para escribir \SI{9.8}{m.s^{-2}} se puede usar \verb+\SI{9.8}{m.s^{-2}}+ o bien \verb+\SI{9.8}{\metre\per\square\second}+.

\section{Paquetes}

Hay determinados paquetes que se cargan de manera autom tica al utilizar la clase \verb+TEMat-article+.
Estos paquetes tienen que estar instalados para que el archivo pueda compilar y se pueda trabajar con las plantillas de la revista.
Por supuesto, no hace falta volver a cargar estos paquetes en el pre mbulo de vuestro documento.
Los paquetes que utilizamos (y que no se pueden quitar) son \verb+expl3+, \verb+xparse+, \verb+l3keys2e+, \verb+silence+, \verb+afterpackage+, \verb+calc+, \verb+etoolbox+, \verb+amsmath+, \verb+amssymb+, \verb+amsthm+, \verb+cabin+, \verb+newtxmath+,  \verb+thmtools+, \verb+babel+ (con las opciones \verb+british+ y \verb+spanish+), \verb+csquotes+, \verb+microtype+, \verb+bm+, \verb+stackengine+, \verb+scalerel+, \verb+siunitx+, \verb+mleftright+, \verb+cancel+, \verb+mathtools+, \verb+graphicx+, \verb+caption+, \verb+subcaption+, \verb+biblatex+ \verb+geometry+, \verb+xcolor+, \verb+enumitem+, \verb+titlesec+, \verb+fancyhdr+, \verb+ccicons+, \verb+xpatch+, \verb+hyperref+, \verb+cleveref+, \verb+fontawesome+ y \verb+environ+.
Adem s, si usas LaTeX est ndar, se cargan tambi n \verb+fontenc+, \verb+inputenc+, \verb+erewhon+ e \verb+inconsolata+ y, si est  disponible, \verb+stix2+; con LuaLaTeX y XeLaTeX se carga \verb+fontspec+, y las tipograf as \verb+Erewhon+ e \verb+Inconsolatazi4+  y, si est  disponible, la tipograf a \verb+STIX2Math+.

Por otra parte, hay otros muchos paquetes que no incluimos obligatoriamente, pero que muchos de vosotros querr is usar.
Para utilizarlos, basta con que los cargu is de la forma habitual.
En general, os recomendamos que no cargu is paquetes que a aden s mbolos matem ticos sin antes comprobar si el s mbolo ya est  disponible.
A modo de ejemplo, el paquete \verb+MnSymbols+ no es compatible con la revista.

\subsection{Figuras, dibujos y diagramas}\label{sec:diagramas}
En muchos art culos puede ser interesante insertar una imagen que ayude a comprender aquello que estamos explicando.
Como norma general, dibujos, diagramas e im genes en general deber an aparecer como figuras autocontenidas.
 Qu  queremos decir con autocontenidas?
Pues que se puedan entender por s  solas, sin necesidad de tener que leer detalladamente el texto.
Para ello, es bueno que est n dentro de un entorno \verb+figure+ y tengan una etiqueta descriptiva usando el comando \verb+caption+.
Pod is ver un ejemplo de ello en la \cref{fig:triangulos}, que adem s tiene subfiguras.
Igual que con las secciones y dem s entornos, se puede usar el comando \verb+\label+ para referirse a ellas m s tarde.

\begin{figure}[hbt!]
  \centering
  \begin{subfigure}[b]{0.31\textwidth}
    \includegraphics[width=\textwidth]{triangulo-precompilado} % el nombre del archivo SIN la extensi n (mejor)
    \caption{Un tri ngulo hecho con TikZ.}\label{fig:triangulo-tikz}
  \end{subfigure}
  \hspace{0.1\textwidth}
  \begin{subfigure}[b]{0.31\textwidth}
    \includegraphics[width=\textwidth]{triangulo-ps-precompilado} % el nombre del archivo SIN la extensi n (mejor)
    \caption{Un tri ngulo hecho con PSTricks.}\label{fig:triangulo-pst}
  \end{subfigure}
  \caption{Ejemplo de figura y subfiguras con TikZ y PSTricks.}\label{fig:triangulos}
\end{figure}

Si os fij is bien en el contenido de las \cref{fig:triangulo-tikz,fig:triangulo-pst}, ver is que hemos usado el comando \verb+\includegraphics+.
Pod is encontrar m s informaci n sobre su uso en la documentaci n del paquete \verb+graphicx+~\cite[sec.~4]{graphics}.

\subsubsection{Informaci n avanzada}\label{sec:diagramas:avanzado}
Hemos probado tanto \verb+tikz+ como \verb+pstricks+.
En caso de usar TikZ, que es el que os recomendamos de los dos, pod is cargar directamente el paquete \verb+tikz+ en el pre mbulo del documento.
Sin embargo, preferimos que aprovech is el paquete \verb+standalone+, y que aprovech is tambi n que hemos preparado un paquete \verb+\TEMat-fonts+ para que pod is generar diagramas y documentos con las tipograf as que usamos en la revista.

La mayor ventaja de usar \verb+standalone+ es que permite compilar las im genes por separado (pod is probarlo con los ejemplos).
Para facilitar la inclusi n de gr ficos con \verb+standalone+, hemos creado dos comandos que se activan al cargar dicho paquete.
Por ejemplo, el archivo \verb+triangulo.tex+ contiene el c digo TikZ para dibujar la \cref{fig:triangulo-tikz}.
Para incluirlo, usar amos el comando \verb+\imagentikz+.
Del mismo modo, \verb+triangulo-ps.tex+ tiene el c digo PSTricks para la \cref{fig:triangulo-pst} y se puede incluir usando \verb+\imagenpstricks+.
Estos comandos aceptan las mismas opciones que \verb+\includegraphics+, pero en este caso no debe incluirse la extensi n del fichero.
As , por ejemplo, la \cref{fig:triangulos} podr amos haberla incluido como sigue:
\begin{verbatim}
\begin{figure}[hbt!]
  \centering
  \begin{subfigure}[b]{0.31\textwidth}
    \imagentikz[width=\textwidth]{triangulo}
    \caption{Un tri ngulo hecho con TikZ.}\label{fig:triangulo-tikz}
  \end{subfigure}
  \hspace{0.1\textwidth}
  \begin{subfigure}[b]{0.31\textwidth}
    \imagenpstricks[width=\textwidth]{triangulo-ps}
    \caption{Un tri ngulo hecho con PSTricks.}\label{fig:triangulo-pst}
  \end{subfigure}
  \caption{Ejemplo de figura y subfiguras con TikZ y PSTricks.}\label{fig:triangulos}
\end{figure}
\end{verbatim}

Es importante que al compilar un documento que use estos comandos, o cualquier otro de los que muestra el paquete \verb+standalone+ en su documentaci n~\cite{standalone} para incluir archivos externos, incluy is la opci n \verb+-shell-escape+ (a veces \verb+-enable-write18+ con MiK\TeX).
Os recomendamos que os le is la \cref{sec:apuntes} para m s informaci n sobre la compilaci n.

Usar TikZ o PSTricks tiene resultados muy bonitos en formato vectorial.
Pero, evidentemente, se pueden incluir figuras hechas de cualquier otra forma, en formato pdf (por ejemplo), de la forma habitual, usando el comando \verb+\includegraphics+ del paquete \verb+graphicx+ que ya hemos incluido autom ticamente.

\begin{nota}
Si ten is conocimientos algo m s avanzados de \LaTeX{}, estos dos comandos para incluir diagramas no son m s que alias del comando \verb+\includestandalone+ con la opci n \verb+mode=buildnew+.
Adem s, en el caso de PSTricks, se debe incluir expl citamente la opci n de compilar con \verb+xelatex+ porque ni \verb+pdflatex+ ni \verb+lualatex+ son compatibles.
\end{nota}

\subsection{Algoritmos}
Pod is usar los paquetes \verb+algpseudocode+ (sin la opci n \verb+noend+) o \verb+program+, aunque tenemos preferencia por \verb+algpseudocode+ (que est  dentro de \verb+algorithmicx+~\cite{algorithmicx}).
Para evitar incompatibilidades, consultadnos antes de usar cualquier otro paquete de algoritmos.

\begin{algoritmo}[escritura a ordenador]
\begin{algorithmic}[1]
\Procedure{Escritura}{}
\If{ordenador apagado}
\State encender ordenador
\If{ordenador apagado}
\If{ordenador enchufado}
\If{no hay corriente}
\State ... \Comment{El lector sabr  qu  hacer}
\Else
\State ordenador roto
\EndIf
\EndIf
\Else
\State ... \Comment{Lo dejamos como ejercicio}
\EndIf
\EndIf
\EndProcedure
\end{algorithmic}
\end{algoritmo}

\subsection{Configuraciones y paquetes prohibidos}
\textbf{Bajo ning n concepto} deben usarse los paquetes o configuraciones que aqu  indicamos:
\begin{itemize}
\item \verb+natbib+ o cualquier otro paquete de bibliograf a.
  Ya usamos \verb+biblatex+ con \verb+biber+, as  que habr a incompatibilidades.
\item \verb+setspace+, \verb+savetrees+ o cualquier paquete que permita cambiar el espaciado del texto.
  Eso implica que aunque carguemos nosotros el paquete \verb+geometry+ no deber ais usar los comandos \verb+geometry+, \verb+newgeometry+ y similares.
  Por supuesto, tampoco deb is usar el comando \verb+\setlength+ para cambiar m rgenes ni ninguna medida del documento, salvo cosas propias vuestras.
\item Cualquier paquete que cambie las tipograf as.
\end{itemize}

\section{Uso de la bibliograf a}

Una de las partes m s importantes de un art culo divulgativo es una buena bibliograf a que permita, por un lado, comprobar de d nde se ha obtenido la informaci n que se ha utilizado para redactar el art culo y, por otro, obtener fuentes en las que leer m s informaci n sobre el tema.
Para crear una bibliograf a adecuada \textbf{se debe utilizar el paquete \texttt{biblatex}} con el \extranj{english}{backend} \verb+biber+ (que cargamos autom ticamente) creando un fichero externo \verb+.bib+ con la bibliograf a completa.
Entre los archivos de la plantilla hay un fichero llamado \verb+plantilla.bib+ que deber a servir como ejemplo de uso de este tipo de archivos.
En cualquier caso, casi siempre se pueden encontrar sitios de los que copiar directamente la bibliograf a en formato bibtex (se pueden buscar las referencias, por ejemplo, en MathSciNet (\url{http://www.ams.org/mathscinet/}) o en Google Scholar (\url{https://scholar.google.com/})).
Para incluir la bibliograf a en el art culo, se debe emplear el comando \verb+\addbibresource{archivobibliografia.bib}+ en el pre mbulo del documento (antes del \verb+\begin{document}+), como se puede ver en el archivo \verb+PlantillaTEMat.tex+, donde \verb+archivobibliografia.bib+ es el nombre que tenga vuestro documento (siempre debe ser un \verb+.bib+).
Para cargar la bibliograf a, se debe escribir \verb+\printbibliography+ al final del documento, donde se quiera que vaya la bibliograf a.

Para referirse a elementos de la bibliograf a, a cada uno de estos elementos se le asigna un nombre en el archivo \verb+.bib+.
Luego basta con usar la instrucci n \verb+\cite{nombrereferencia}+ para que aparezca, proporcionando adem s un v nculo a la bibliograf a del final del art culo.

Toda la informaci n sobre esto, y otras muchas cosas relacionadas con la bibliograf a, la pod is encontrar en la documentaci n de \verb+biblatex+~\cite{biblatex}.
Si quer is buscar la versi n m s actualizada, es f cil encontrarla en su p gina de CTAN.\footnote{Acr nimo de \emph{Comprehensive \protect\TeX{} Archive Network}, ver \url{https://www.ctan.org/pkg/biblatex} y \url{https://www.ctan.org/}}
En esa misma web se pueden encontrar manuales de uso de casi todos los paquetes de \LaTeX.

\begin{nota}
El comando \verb+\printbibliography+ har  que aparezcan solo los elementos que se han citado a lo largo del art culo.
Si hay alguna referencia que no se ha citado pero que tambi n deber a aparecer, se debe usar el comando \verb+\nocite{nombrereferencia}+.
Para que aparezcan todas las que no se han citado, se puede usar un asterisco: \verb+\nocite{*}+.
\end{nota}

\begin{nota}
Para mejorar la calidad del texto, es recomendable no usar la salida del comando \verb+\cite+ como un sustantivo.
Como ejemplo, es mejor decir <<en la documentaci n de \texttt{biblatex}~\cite{biblatex} hay m s informaci n>> que <<en~\cite{biblatex} hay m s informaci n>>.

Tambi n existe el comando \verb+\citet+ para referirse a un trabajo nombrando al autor.
Es bueno cuando quieres decir cosas del tipo <<\citet{bezosOM} muestra gran cantidad de informaci n sobre c mo se deber an escribir las matem ticas a ordenador>>. Otra opci n es usar \verb+\citeauthor+ combinado con \verb+\cite+, como en <<\citeauthor{bezosOM} muestra gran cantidad de informaci n sobre c mo se deber an escribir las matem ticas a ordenador~\cite{bezosOM}>>, aunque recomendamos usar \verb+\citet+.
\end{nota}

\begin{nota}
Si ten is problemas con la bibliograf a, pod is intentar a adir la opci n \verb+bibtex+ al cargar la clase (por ejemplo, \verb+\documentclass[bibtex,anon]{TEMat-article}+).
Esta opci n cambia la configuraci n para que se use el \extranj{english}{backend} \verb+bibtex+ en lugar de \verb+biber+ a la hora de extraer los datos bibliogr ficos del archivo con extensi n \verb+bib+.

Esta es una opci n que hemos a adido por compatibilidad para instalaciones m s antiguas, pero recomendamos que no se use.
Las versiones finales de los art culos aceptados, cuando las generemos para la revista, no podr n usar esta opci n.
\end{nota}

\section{Algunos apuntes}\label{sec:apuntes}
La compilaci n de documentos escritos con \LaTeX{} es un poco odisea, los que se is usuarios habituales ya lo sabr is.
Si no se usan paquetes como \verb+minted+ o lo expuesto en la \cref{sec:diagramas:avanzado}, basta con ejecutar los comandos
\begin{verbatim}
pdflatex PlantillaTEMat
biber PlantillaTEMat
pdflatex PlantillaTEMat
pdflatex PlantillaTEMat
\end{verbatim}
desde la consola, sustituyendo el comando \verb+pdflatex+ por \verb+lualatex+ si es el caso y \verb+PlantillaTEMat+ por el nombre del archivo \TeX{} que se est  compilando.

Si se usan paquetes que tienen que hacer llamadas externas (\verb+minted+ o im genes TikZ o PSTricks externalizadas) es necesario compilar en modo \extranj{english}{shell escape}:
\begin{verbatim}
pdflatex -shell-escape PlantillaTEMat
biber PlantillaTEMat
pdflatex -shell-escape PlantillaTEMat
pdflatex -shell-escape PlantillaTEMat
\end{verbatim}

Por supuesto, si hab is pasado la opci n \verb+bibtex+ a la clase, la segunda orden deber is sustituirla por una llamada a \verb+bibtex+.

Una forma m s f cil de compilar el documento es usar la herramienta \verb+latexmk+, que viene por defecto en las  ltimas versiones de \TeX~Live y de MiK\TeX{}, aunque en Windows es necesario instalar alguna implementaci n de Perl como ActivePerl\footnote{\url{http://www.activestate.com/activeperl/downloads}} para que funcione.
Con esto, y mediante el archivo \verb+latexmkrc+ que incluimos con la plantilla, basta con hacer
\begin{verbatim}
latexmk PlantillaTEMat
\end{verbatim}
y el propio \verb+latexmk+ se encarga de llamar todas las veces que haga falta a cada una de las herramientas.
Se puede editar el archivo \verb+latexmk+ con un editor de texto plano para cambiar ciertas opciones.

Todav a m s f cil es usar Overleaf\footnote{\url{https://www.overleaf.com/}}.
Basta con subir los archivos de la plantilla y \emph{voil }, funciona.
Pod is configurarlo para que use LuaLaTeX, ya que es el formato con el que compilaremos nosotros.
Es cierto que pod is usar pdfLaTeX si lo prefer s, la apariencia ser  pr cticamente igual pero compilar  m s r pido.

Si us is un editor con botones para compilar, buscad la informaci n para configurarlo.
Muchas veces permiten el uso de \verb+latexmk+, as  que lo mejor ser a hacer que as  fuera.
Si vais a compilar la plantilla de ejemplo y us is \verb+minted+ o alguno de los comandos para incluir im genes TikZ o PSTricks, aseguraos de que el comando de compilaci n tenga activado el \extranj{english}{shell escape}.


\section{Usos avanzados}

La clase \verb+TEMat-article+ incorpora opciones para usos m s avanzados as  como soporte para otros idiomas distintos al castellano de forma preliminar.
M s adelante actualizaremos esta documentaci n para explicar c mo se accede a estas opciones que, por otro lado, en general no estar n autorizadas para su uso en TEMat.


\section{Conclusiones}

En funci n del art culo, puede ser adecuado incluir para terminar una secci n de conclusiones, preguntas abiertas sobre el tema, formas de obtener m s informaci n...
Algo que permita cerrar el art culo.

Esta plantilla no puede ser, ni mucho menos, una introducci n al enorme mundo de \LaTeX.
Para aprender m s, recomendamos buscar informaci n por internet, donde hay gran cantidad de fuentes~\cite{wiki:latex, stackexchange, texample}.

Si necesit is algo o ten is dudas sobre la revista, escribid a \email{temat@temat.es}.

\nocite{*}
\printbibliography[heading=bibintoc]
\end{document}